\subsection{Bidirectional NLI examples}\label{app:nli-examples}

For each of mean bidirectional NLI $\in \{0.5, 0.7, 0.9\}$, we show one (original, rewrite) pair generated by the na\"ive rewriting method. Mean NLI is $\tfrac{1}{2}(\mathrm{NLI}_{\rightarrow} + \mathrm{NLI}_{\leftarrow})$ where $\rightarrow$ is the probability that the rewrite entails the original (forward), and $\leftarrow$ the reverse (backward). Embedding cosine similarity is computed with \texttt{intfloat/e5-large-v2}.

\subsubsection*{Mean NLI $\approx$ 0.5}

\begin{description}\setlength{\itemsep}{0pt}\setlength{\parsep}{0pt}
  \item[Proposition] The UK should legally recognize and protect all forms of transgender identity.
  \item[Method] na\"ive~~(rewriter: Qwen2.5-1.5B)
  \item[NLI fwd / bwd / mean] 0.064 / 0.930 / 0.497
  \item[Embedding cosine] 0.880
\end{description}

\noindent\textbf{Original.}~~\textit{We live in a very difficult and troubled world where religion and identity seem to be the centre of a lot of issues. I feel that is we could iron out the opinions that transgender people are less entitled to an opinion, then society in general may be able to live side by side with all people no matter what their identity is. Transgender people have as much right as anyone else of their opinions and I feel sure many people could learn a lot of they could accept this and listen to them. It must be very frustrating for transgender people, which may not bring out the best in them.}

\noindent\textbf{Rewrite.}~~\textit{In a world grappling with complex issues like religion and identity, it is crucial to recognize and protect the rights of all individuals, regardless of their transgender identity. By fostering an environment where all people, including transgender individuals, feel valued and respected, society can progress together. Transgender individuals, like others, deserve the right to express their opinions and should be encouraged to do so. Their experiences and perspectives can offer enriching insights that can help bridge gaps and promote understanding. It is essential that society listens to and embraces these voices, acknowledging the struggles and challenges they face. This approach can lead to a more inclusive and harmonious society where everyone can thrive.}

\subsubsection*{Mean NLI $\approx$ 0.7}

\begin{description}\setlength{\itemsep}{0pt}\setlength{\parsep}{0pt}
  \item[Proposition] The UK should mandate COVID-19 vaccinations for all eligible citizens.
  \item[Method] na\"ive~~(rewriter: Qwen2.5-1.5B)
  \item[NLI fwd / bwd / mean] 0.405 / 0.991 / 0.698
  \item[Embedding cosine] 0.908
\end{description}

\noindent\textbf{Original.}~~\textit{We saw so much delay during Covid and who knows how many lives were lost due to indecision, conspiracies and so on. I remember a fitness instructor who didn't want the vaccine because he had believed some false information about the vaccine and he subsequently died. He almost certainly would have lived if he had had the vaccine. There is a responsibility on ourselves to act in the best interests of the whole public, and we saw too many people acting selfishly. Therefore let's take away the indecision and move more quickly by making it mandatory to be vaccinated. It would save lives.}

\noindent\textbf{Rewrite.}~~\textit{''''''A delay during the pandemic exposed our collective hesitancy. An example is a fitness instructor who died after rejecting the vaccine due to false beliefs. Had he been vaccinated, he likely would have survived. It's our responsibility to prioritize public welfare. We observed too much selfishness. Mandating vaccinations is necessary to save lives. It would take away indecision and expedite the process. Let's move forward with mandatory vaccinations.''''''}

\subsubsection*{Mean NLI $\approx$ 0.9}

\begin{description}\setlength{\itemsep}{0pt}\setlength{\parsep}{0pt}
  \item[Proposition] The UK should ban internet use for under-14-year-olds.
  \item[Method] na\"ive~~(rewriter: Qwen2.5-1.5B)
  \item[NLI fwd / bwd / mean] 0.922 / 0.883 / 0.902
  \item[Embedding cosine] 0.891
\end{description}

\noindent\textbf{Original.}~~\textit{The Internet is growing part of learning and development for younger people. Taking this away could hinder and not progress their education like it should. It could cause more stress leading to higher levels of school drop out and under achievement. There is a lot of reliance on technology and it is a big part of what people do and how people communicate. Taking this away would make males feel isolated and unaware of what others may be going on about. It could end up almost like a giant social click of boys against girls. If the concern is due to what can be found online then the ban would need to be implemented across the board not just aimed at males. There are resources which make it easy to monitor and restrict what can and cannot be accessed which may be more beneficial}

\noindent\textbf{Rewrite.}~~\textit{The Internet plays a crucial role in younger people's education and development. Restricting internet use for underage individuals could stifle their learning and hinder their progress. This measure could also lead to increased stress and higher rates of school dropout and underachievement. Reliance on technology is pervasive; removing it would isolate males and limit their understanding of others. Implementing a blanket ban targeting both genders would make boys feel isolated and disconnected from their peers. If concerns about online content are valid, a comprehensive ban across all platforms would be more effective than restricting access to males alone.}

\backlink{app:nli-examples}
